\chapter{Extensions du cadre d'inventaire}\label{ch:hf:extensions-of-the-inventory-framework}

La solution exacte du chapitre~3 conseille à un teneur de marché qui détient quatre unités de les offrir sous le mid. Les teneurs de marché réels cessent plutôt d'acheter~:
ils bornent l'inventaire. Ils détiennent aussi plusieurs carnets corrélés à la fois, ils ont une idée de la direction du prix, et ils savent
que la prochaine transaction a davantage de chances de venir de quelqu'un qui en sait plus qu'eux. Chacun de ces éléments entre dans le même système linéaire. Dans
les simulations de ce chapitre, coter deux actifs dont les prix sont corrélés à 0.8 comme un seul carnet au lieu de deux réduit de 21~\% le risque d'inventaire pour le même
profit espéré~; un teneur de marché qui connaît la dérive tient la position que vise la formule~; et celui dont les exécutions sont suivies d'un mouvement adverse
a intérêt à élargir chaque cotation de la moitié de ce mouvement, et non de la totalité.

\section{Bornes d'inventaire et solution asymptotique}

\begin{definition}[Borne d'inventaire]\label{def:hf:extensions-of-the-inventory-framework:bound}
Une \emph{borne d'inventaire}\index{borne d'inventaire} $\bar q$ est une limite stricte sur la position du teneur de marché~: en $q=\bar q$ il ne cote pas à l'achat, en
$q=-\bar q$ il ne cote pas à la vente, de sorte que $|q_t|\le\bar q$ à tout instant.
\end{definition}

La borne fait partie de l'espace des états de la \cref{prop:hf:inventory-models:cj} du chapitre~3, ce qui explique que cette solution soit exacte sur une grille finie. Son
premier usage est pratique~: quel que soit le modèle, la position ne peut pas s'emballer. Le second est analytique~: sur une grille finie, le comportement
à long horizon de la solution est un problème de vecteur propre.

\begin{proposition}[Les cotations stationnaires]\label{prop:hf:extensions-of-the-inventory-framework:stationary}
Lorsque $T-t\to\infty$, les profondeurs du modèle à pénalité courante du chapitre~3 convergent vers des profondeurs indépendantes du temps données par les mêmes formules avec
$h=\tfrac1k\ln v$, où $v$ est le vecteur propre associé à la plus grande valeur propre de $M$.
\end{proposition}

\begin{proof}
$M$ est symétrique, avec des coefficients hors diagonale positifs juste à côté de la diagonale~; sa plus grande valeur propre $\mu_1$ est donc simple, et son vecteur propre $v$ peut
être pris positif (Perron--Frobenius). Alors $\omega(t)=e^{M(T-t)}z=\sum_ie^{\mu_i(T-t)}(v_i^\top z)v_i$ est dominé par le premier terme lorsque $T-t\to\infty$,
puisque $v^\top z>0$. Les profondeurs ne dépendent de $\omega$ que par les rapports selon $q$, de sorte que le facteur $e^{\mu_1(T-t)}(v^\top z)$ disparaît.
\end{proof}

Avec les paramètres du chapitre~3 ($A=140$, $k=1.5$, $\phi=0.5$), un horizon d'une unité de temps est déjà long~: les profondeurs à son début coïncident avec les
profondeurs stationnaires à huit décimales près. Guéant, Lehalle et Fernandez-Tapia ont résolu de la même façon la version à utilité exponentielle et ont approché
le vecteur propre sous forme fermée. La même approximation pour la pénalité courante donne une formule assez simple pour figurer dans une fonction de prix.

\begin{proposition}[Cotations stationnaires en forme fermée]\label{prop:hf:extensions-of-the-inventory-framework:closed}
Si le vecteur propre est approximativement gaussien, $v_q\propto e^{-\beta q^2/2}$ avec $\beta$ petit, alors $\beta^2=\phi ke/A$ et, avec
$c=\beta/k=\sqrt{\phi e/(Ak)}$,
\[
\delta^b(q)\simeq\frac1k+\frac{2q+1}{2}\,c,\qquad \delta^a(q)\simeq\frac1k-\frac{2q-1}{2}\,c .
\]
\end{proposition}

\begin{proof}
L'équation aux valeurs propres s'écrit $Ae^{-1}(v_{q-1}+v_{q+1})-\phi kq^2v_q=\mu_1v_q$. Avec l'ansatz gaussien, $v_{q\pm1}/v_q=e^{-\beta(\pm2q+1)/2}$, dont la somme
vaut $2e^{-\beta/2}\cosh(\beta q)\simeq2-\beta+\beta^2q^2$ à l'ordre deux. L'identification des termes en $q^2$ donne $Ae^{-1}\beta^2=\phi k$. Alors
$h(q)-h(q\pm1)=\tfrac1k\ln(v_q/v_{q\pm1})=\tfrac{\beta}{2k}(\pm2q+1)$.
\end{proof}

Chaque unité d'inventaire décale les deux cotations de $c$, et la fourchette reste égale à $2/k+c$. Pour $\gamma\to0$ avec $\phi=\tfrac12\gamma\sigma^2$, la
formule de Guéant--Lehalle--Fernandez-Tapia se réduit à celle-ci. L'approximation exige que $\beta$ soit petit, c'est-à-dire que la pénalité soit petite devant le taux d'exécution~:

\begin{center}
\small
\begin{tabular}{@{}lccccc@{}}
\toprule
profondeur à l'achat, $\phi=0.5$ & $q=0$ & $q=1$ & $q=2$ & $q=3$ & écart maximal, $|q|\le10$\\
\midrule
exacte, stationnaire & 0.708 & 0.790 & 0.872 & 0.953 & \\
forme fermée & 0.707 & 0.787 & 0.868 & 0.948 & 0.034\\
\midrule
écart maximal pour $\phi=5$ / $\phi=50$ & & & & & 0.70 / 5.0\\
\bottomrule
\end{tabular}
\end{center}

La \Cref{fig:hf:extensions-of-the-inventory-framework:closed} montre la défaillance pour $\phi=5$~: la profondeur exacte se courbe pour les grands inventaires, là où le
vecteur propre est loin d'être gaussien~; la droite ne le fait pas.

\begin{omfigure}
\begin{tikzpicture}
\begin{axis}[width=11cm,height=5.2cm,xlabel={\small inventaire $q$},ylabel={\small profondeur à l'achat},tick label style={font=\footnotesize},
  xmin=-8,xmax=8,grid=major,grid style={gray!25},table/col sep=comma,legend style={font=\scriptsize,at={(0.5,-0.3)},anchor=north},legend columns=4]
\addplot[omDef,thick] table[x=q,y=exact05]{figdata/hft/04-extensions-of-the-inventory-framework/closed.csv};
\addplot[omDef,thick,dashed] table[x=q,y=closed05]{figdata/hft/04-extensions-of-the-inventory-framework/closed.csv};
\addplot[omThm,thick] table[x=q,y=exact5]{figdata/hft/04-extensions-of-the-inventory-framework/closed.csv};
\addplot[omThm,thick,dashed] table[x=q,y=closed5]{figdata/hft/04-extensions-of-the-inventory-framework/closed.csv};
\legend{{exacte, $\phi=0.5$},{forme fermée, $\phi=0.5$},{exacte, $\phi=5$},{forme fermée, $\phi=5$}}
\end{axis}
\end{tikzpicture}

\omcaption{Profondeur stationnaire à l'achat en fonction de l'inventaire~: exacte (vecteur propre) et en forme fermée, pour deux pénalités courantes ($A=140$, $k=1.5$, inventaire
dans $\pm20$). Données~: \texttt{firm.multimm}.}\label{fig:hf:extensions-of-the-inventory-framework:closed}
\end{omfigure}

\section{Multiples actifs et inventaire corrélé}

\begin{definition}[Market making multi-actifs]\label{def:hf:extensions-of-the-inventory-framework:multi}
Le \emph{market making multi-actifs}\index{market making multi-actifs} cote plusieurs instruments à partir d'un vecteur d'inventaire $q\in\mathbb Z^d$ et d'une
charge de risque sur l'ensemble du carnet, $\phi\,q^\top\Sigma q$ par unité de temps, où $\Sigma$ est la covariance des variations de prix des instruments, de sorte que chaque
cotation dépend de l'inventaire de chaque instrument.
\end{definition}

\begin{proposition}[Cotations exactes multi-actifs]\label{prop:hf:extensions-of-the-inventory-framework:multi}
Si les instruments partagent $A$ et $k$, la substitution $h=\tfrac1k\ln\omega$ transforme le problème sur la grille $\{-\bar q,\dots,\bar q\}^d$ en le
système linéaire $\partial_t\omega+M\omega=0$, avec $M_{q,q}=-\phi k\,q^\top\Sigma q$ et $Ae^{-1}$ entre les inventaires voisins $q\pm e_i$~; les
profondeurs de l'instrument $i$ sont $\delta^{b,a}_i(q)=\tfrac1k+h(q)-h(q\pm e_i)$, et la \cref{prop:hf:extensions-of-the-inventory-framework:stationary}
reste valable sur la grille.
\end{proposition}

\begin{proof}
Comme au chapitre~3, instrument par instrument~: le supremum de chaque côté apporte
\[
\tfrac1k\,Ae^{-1}\,e^{k(h(q\pm e_i)-h(q))},
\]
qui est linéaire en $\omega$ après la substitution, car chaque instrument a le même $k$.
\end{proof}

Le $k$ commun est le prix de l'exactitude~; avec des $k_i$ différents, le système n'est plus linéaire et la littérature consacrée au \omterm{def:hf:extensions-of-the-inventory-framework:multi}{market making multi-actifs} (Guéant 2017) recourt à des
approximations quadratiques de la fonction valeur. La grille compte $(2\bar q+1)^d$ points, soit 289 pour deux instruments avec $\bar q=8$~: acceptable pour une paire, mais
pas pour une centaine de noms.

La corrélation crée un skew croisé. Sous la politique conjointe (la \cref{fig:hf:extensions-of-the-inventory-framework:cross}), un teneur de marché court de quatre
unités du second instrument cote à l'achat sur le premier à 0.572 du mid lorsqu'il est à plat dans celui-ci, contre 0.700 lorsqu'il est à plat dans les deux~: une position longue dans le
premier couvre la position courte dans le second. Instrument par instrument, le cours acheteur serait de 0.700 quelle que soit l'autre position.

\begin{omfigure}
\begin{tikzpicture}
\begin{axis}[width=11cm,height=5.2cm,xlabel={\small inventaire de l'instrument 1},ylabel={\small profondeur acheteuse de l'instrument 1},tick label style={font=\footnotesize},
  xmin=-6,xmax=6,grid=major,grid style={gray!25},table/col sep=comma,legend style={font=\scriptsize,at={(0.5,-0.3)},anchor=north},legend columns=3]
\addplot[omDef,thick,mark=*,mark size=1pt] table[x=q1,y=short4]{figdata/hft/04-extensions-of-the-inventory-framework/cross.csv};
\addplot[black,thick,mark=square*,mark size=1pt] table[x=q1,y=flat]{figdata/hft/04-extensions-of-the-inventory-framework/cross.csv};
\addplot[omThm,thick,mark=triangle*,mark size=1.3pt] table[x=q1,y=long4]{figdata/hft/04-extensions-of-the-inventory-framework/cross.csv};
\legend{{position courte de 4 sur l'instrument 2},{à plat sur l'instrument 2},{position longue de 4 sur l'instrument 2}}
\end{axis}
\end{tikzpicture}

\omcaption{La profondeur acheteuse de la politique conjointe pour l'instrument~1 en fonction de son propre inventaire, pour trois inventaires de l'instrument~2 (deux actifs
de corrélation des prix 0.8, pénalité $\phi=0.1$ sur $q^\top\Sigma q$, $A=140$, $k=1.5$, $\sigma=2$). Données~: \texttt{hf\_extensions.cross\_skew}.}\label{fig:hf:extensions-of-the-inventory-framework:cross}
\end{omfigure}

La politique à skew croisé rapporte. Sur 1\,000 trajectoires simulées de cinq unités de temps, avec les deux mids corrélés à 0.8, la politique conjointe et la
politique instrument par instrument (la même pénalité appliquée à la variance propre de chaque actif) rapportent presque autant, mais le livre conjoint porte moins de risque d'inventaire, mesuré comme la
moyenne temporelle de $q^\top\Sigma q$ (la \cref{fig:hf:extensions-of-the-inventory-framework:frontier}). Au gain moyen de la politique conjointe pour $\phi=0.1$,
673.3, la frontière instrument par instrument porte 30.8 unités de risque d'inventaire et la politique conjointe 24.3~: 21~\% de moins. Sur les six pénalités, la
réduction va de 10~\% à 22~\%.

\begin{omfigure}
\begin{tikzpicture}
\begin{axis}[width=11cm,height=5.2cm,xlabel={\small risque d'inventaire (moyenne temporelle de $q^\top\Sigma q$)},ylabel={\small gain moyen},
  tick label style={font=\footnotesize},grid=major,grid style={gray!25},table/col sep=comma,
  legend style={font=\scriptsize,at={(0.5,-0.3)},anchor=north},legend columns=2]
\addplot[omThm,thick,mark=*,mark size=1.3pt] table[x=risk_joint,y=mean_joint]{figdata/hft/04-extensions-of-the-inventory-framework/frontier2.csv};
\addplot[omDef,thick,dashed,mark=square*,mark size=1.3pt] table[x=risk_sep,y=mean_sep]{figdata/hft/04-extensions-of-the-inventory-framework/frontier2.csv};
\legend{livre conjoint,instrument par instrument}
\end{axis}
\end{tikzpicture}

\omcaption{Gain moyen en fonction du risque d'inventaire pour deux actifs corrélés à 0.8, cotés en un seul livre ou instrument par instrument, pour des pénalités $\phi$ allant de 0.01
(à droite) à 0.5 (à gauche)~; 1\,000 trajectoires communes de cinq unités de temps. Données~: \texttt{hf\_extensions.frontier2}.}\label{fig:hf:extensions-of-the-inventory-framework:frontier}
\end{omfigure}

\section{Signaux et dérive au sein du problème de contrôle}

\begin{definition}[Cotation ajustée de la dérive]\label{def:hf:extensions-of-the-inventory-framework:drift}
La \emph{cotation ajustée de la dérive}\index{cotation ajustée de la dérive} (drift-adjusted quoting) résout le problème de contrôle du teneur de marché avec la dérive prévisionnelle $\mu$ du mid
en l'intégrant, de sorte que la prévision déplace les cotations par la valeur de la détention d'inventaire plutôt que de s'y ajouter comme un skew distinct.
\end{definition}

Avec $dS_t=\mu\,dt+\sigma\,dW_t$, détenir $q$ rapporte $\mu q$ par unité de temps, et la récompense instantanée vaut $\mu q-\phi q^2=-\phi(q-q_0)^2+\text{constante}$ avec
$q_0=\mu/(2\phi)$.

\begin{proposition}[Une dérive fixe une position cible]\label{prop:hf:extensions-of-the-inventory-framework:target}
Avec une dérive constante $\mu$, la diagonale de $M$ devient $k(\mu q-\phi q^2)$, et dans la forme fermée de la
\cref{prop:hf:extensions-of-the-inventory-framework:closed}, le teneur de marché cote comme si son inventaire était $q-q_0$~: il vise la position
$q_0=\mu/(2\phi)$.
\end{proposition}

\begin{proof}
La récompense instantanée entre dans la diagonale de $M$ multipliée par $k$, comme $-\phi q^2$ auparavant. En complétant le carré, l'équation aux valeurs propres de la
\cref{prop:hf:extensions-of-the-inventory-framework:closed} reste vraie avec $q$ remplacé par $q-q_0$, à une constante près absorbée dans $\mu_1$.
\end{proof}

La formule de Guéant, Lehalle et Fernandez-Tapia avec dérive a la même forme~: un décalage de l'inventaire de $\mu/(\gamma\sigma^2)$, qui vaut $q_0$ pour
$\phi=\tfrac12\gamma\sigma^2$. Cartea et Wang introduisent un signal alpha stochastique dans le problème, dans le même esprit. Dans la simulation avec $\mu=2$,
$\phi=0.5$ (donc $q_0=2$), le teneur de marché ajusté à la dérive détient en moyenne 1.97 unité et gagne 355.6, contre 335.7 pour le teneur de marché qui ignore la
dérive et reste à plat~: 5.9~\% de plus, pour un écart type de 18.0 au lieu de 15.8. La prévision ne vaut de l'argent que par la position qu'elle fait
détenir au teneur de marché~; le chapitre~7 construit de telles prévisions à l'horizon de quelques secondes.

\section{Sélection adverse au sein du modèle}

La sélection adverse entre dans le modèle de deux façons, et elles appellent des élargissements différents.

\begin{proposition}[Élargir du coût, ou de la moitié du mouvement]\label{prop:hf:extensions-of-the-inventory-framework:adverse}
(i) Si chaque exécution coûte en moyenne $\varepsilon$ (le prix d'exécution est en moyenne inférieur de $\varepsilon$ au mid), le terme hors diagonale de $M$ devient
$Ae^{-1-k\varepsilon}$ et les deux profondeurs optimales augmentent de $\varepsilon$ à chaque inventaire. (ii) Si au contraire le mid se déplace de façon permanente de $\xi$ contre
le teneur de marché après chacune de ses exécutions, les profondeurs optimales sont celles de (i) avec $\varepsilon=\xi/2$.
\end{proposition}

\begin{proof}
(i) Chaque supremum devient $\sup_\delta Ae^{-k\delta}(\delta-\varepsilon+\Delta h)$~: l'optimum vaut $1/k+\varepsilon-\Delta h$ et sa valeur porte un facteur
$e^{-k\varepsilon}$. (ii) Après une exécution à l'achat, toute la position $q+1$ est réévaluée de $-\xi$, après une exécution à la vente $q-1$ de $+\xi$~: les crochets valent
$\delta-\xi(q+1)+h(q+1)-h(q)$ et $\delta+\xi(q-1)+h(q-1)-h(q)$. Avec $h(q)=\tilde h(q)+\tfrac12\xi q^2$, les deux deviennent $\delta-\tfrac12\xi+\Delta\tilde
h$, soit le problème (i) pour $\tilde h$ avec $\varepsilon=\xi/2$, et les profondeurs exprimées à l'aide de $\tilde h$ sont celles de (i).
\end{proof}

La moitié, parce que le mouvement défavorable à l'unité qui vient d'être exécutée réévalue aussi le reste de la position, et qu'en moyenne la position que laisse l'exécution est du
bon côté du mouvement aussi souvent que du mauvais. La simulation de (ii) avec $\xi=0.3$, $\phi=0.5$ le confirme, à condition que le pas de temps soit
assez fin~:

\begin{center}
\small
\begin{tabular}{@{}lccc@{}}
\toprule
élargissement de chaque profondeur & aucun & $\xi/2=0.15$ & $\xi=0.3$\\
\midrule
gain moyen & 262.7 & 267.2 & 258.9\\
gain diminué de la pénalité & 253.1 & 258.6 & 251.2\\
exécutions & 498 & 396 & 315\\
gain diminué de la pénalité, pas de temps quatre fois plus grossier & 270.8 & 270.2 & 259.3\\
\bottomrule
\end{tabular}
\end{center}

Avec le pas grossier (0.005, lorsqu'une cotation à plat a environ 24~\% de chance d'être exécutée à chaque pas), la politique non ajustée paraissait marginalement la meilleure~; divisé par quatre, le pas
reproduit la proposition. La règle du livre~7 s'applique~: une simulation sur laquelle s'appuie un résultat nommé est relancée avec un quart du
pas de temps.

\section{Tutoriel~: un seul système linéaire, quatre extensions}\label{tut:hf:extensions-of-the-inventory-framework:tut}

\begin{tutorial}
\textbf{Objectif.} Calculer des cotations stationnaires, en forme fermée, multi-actifs, ajustées à la dérive et ajustées à la sélection adverse à partir d'une même construction, et tester
chacune en simulation. \textbf{État final~:} les trois tableaux et les trois figures du chapitre.
\begin{tutsteps}
\item \textbf{Stationnaire et forme fermée.} Le vecteur propre de Perron de $M$ donne les profondeurs de long terme~; l'approximation gaussienne donne la
formule (le \cref{lst:hf:extensions-of-the-inventory-framework:stationary}).
\omcode{code/firm/multimm/firm_multimm.py}{49}{69}{Profondeurs stationnaires à partir du vecteur propre, et leur forme fermée.}\label{lst:hf:extensions-of-the-inventory-framework:stationary}
\item \textbf{Deux actifs.} La même construction sur la grille des paires d'inventaires
(le \cref{lst:hf:extensions-of-the-inventory-framework:multi})~; \texttt{simulate\_multi} tire des mids corrélés.
\omcode{code/firm/multimm/firm_multimm.py}{76}{94}{Le système linéaire multi-actifs et son vecteur propre.}\label{lst:hf:extensions-of-the-inventory-framework:multi}
\item \textbf{Frontières.} \texttt{hf\_extensions.frontier2()} lance les deux politiques pour six pénalités sur 1\,000 trajectoires communes~;
\texttt{equal\_capture()} interpole la frontière instrument par instrument à chaque gain moyen conjoint.
\item \textbf{Dérive et sélection adverse.} \texttt{drift\_case(adjust)} et \texttt{adverse\_case(adjust, dt)}.
\end{tutsteps}
\textbf{Pour aller plus loin.} Donner aux deux actifs des taux d'exécution différents et comparer la solution sur grille avec une approximation quadratique~; faire changer de signe la
dérive au hasard (un signal de Markov) et l'ajouter à l'état.
\end{tutorial}

\section{Construction~: les extensions multi-actifs et à signal}\label{bld:hf:extensions-of-the-inventory-framework:multimm}

\begin{build}
\bfield{Rôle} Cotations à long horizon qu'un quoteur de production peut évaluer en microsecondes (une table ou une formule), pour un ou plusieurs instruments,
avec une dérive et une sélection adverse à l'intérieur du problème.
\bfield{Interface} \texttt{firm.multimm}~: \texttt{stationary(A, k, phi, qmax, mu, eps)}, \texttt{asymptotic(A, k, phi, q, mu, eps)},
\texttt{MultiStationary(A, k, phi, cov, qmax).depths(q)}, \texttt{simulate\_multi(policy, T, dt, cov, A, k, paths, seed, qmax)},
\texttt{table\_policy}~; \texttt{firm.invmm.CJSolution(\dots, mu, eps)} pour les horizons finis.
\bfield{Règles} La \omterm{def:hf:extensions-of-the-inventory-framework:bound}{borne d'inventaire} n'est jamais franchie. Un déplacement permanent après exécution $\xi$ est transmis comme $\varepsilon=\xi/2$. Les simulations qui étayent un
résultat imprimé sont répétées avec un quart du pas de temps.
\bfield{Tests d'acceptation} \texttt{code/firm/multimm/tests/}~: les profondeurs stationnaires égalent les profondeurs à horizon fini au début d'un horizon long~; la
forme fermée à 0.05 près des profondeurs exactes pour une pénalité faible et son skew par unité égal à $c$~; une dérive décale les cotations de $q_0$~; $\varepsilon$
décale les deux profondeurs en forme fermée de $\varepsilon$~; deux actifs non corrélés se cotent comme deux actifs seuls~; avec corrélation, une position courte dans l'un
des actifs abaisse la profondeur acheteuse de l'autre~; le simulateur est déterministe et borné.
\bfield{Pour aller plus loin} Des $k_i$ différents avec une approximation quadratique de $h$~; une dérive modulée par un processus de Markov dans l'état.
\end{build}

\begin{omsources}
\item O.~Guéant, C.-A.~Lehalle et J.~Fernandez-Tapia, «~Dealing with the inventory risk~: a solution to the market making problem~»,
\emph{Mathematics and Financial Economics} 7(4), 2013.
\item O.~Guéant, «~Optimal market making~», \emph{Applied Mathematical Finance} 24(2), 2017.
\item Á.~Cartea et Y.~Wang, «~Market making with alpha signals~», \emph{International Journal of Theoretical and Applied Finance} 23(3), 2020.
\end{omsources}

\section{Exercices}

\begin{exercise}[$\star$]\label{exo:hf:extensions-of-the-inventory-framework:1}
Calculer $c=\sqrt{\phi e/(Ak)}$ pour $A=140$, $k=1.5$, $\phi=0.5$, et les profondeurs acheteuse et vendeuse en forme fermée à $q=2$.
\end{exercise}

\begin{exercise}[$\star$]\label{exo:hf:extensions-of-the-inventory-framework:2}
Un teneur de marché avec $\phi=0.5$ anticipe une dérive de $\mu=1$. Déterminer la position qu'il vise.
\end{exercise}

\begin{exercise}[$\star$]\label{exo:hf:extensions-of-the-inventory-framework:3}
Déterminer le nombre de points de grille nécessaires à la solution exacte multi-actifs pour trois instruments avec $\bar q=8$, puis pour dix.
\end{exercise}

\begin{exercise}[$\star\star$]\label{exo:hf:extensions-of-the-inventory-framework:4}
Montrer que la fourchette en forme fermée $\delta^a+\delta^b$ ne dépend pas de l'inventaire, et la calculer pour les paramètres de l'exercice 1.
\end{exercise}

\begin{exercise}[$\star\star$]\label{exo:hf:extensions-of-the-inventory-framework:5}
Expliquer pourquoi la politique conjointe cote plus agressivement à l'achat pour l'actif 1 lorsqu'elle est courte sur l'actif 2, et dire si elle le ferait encore si la corrélation était négative.
\end{exercise}

\begin{exercise}[$\star\star$]\label{exo:hf:extensions-of-the-inventory-framework:6}
Expliquer en mots pourquoi un déplacement permanent après exécution $\xi$ appelle un élargissement de $\xi/2$ plutôt que de $\xi$.
\end{exercise}

\begin{exercise}[$\star\star\star$]\label{exo:hf:extensions-of-the-inventory-framework:7}
\emph{Programmation.} Avec \texttt{stationary(140, 1.5, 0.5, 20, mu=2.0)}, donner les profondeurs acheteuse et vendeuse en $q=0$ et en $q=2$, et les comparer aux
profondeurs en $q=0$ sans dérive.
\end{exercise}

\begin{exercise}[$\star\star\star$]\label{exo:hf:extensions-of-the-inventory-framework:8}
\emph{Trouver la faille.} «~Notre simulation montre que l'élargissement pour sélection adverse fait perdre de l'argent, donc on cote à la profondeur de monopole et on l'ignore.~»
\end{exercise}

\section{Problème~: Deux portefeuilles, un seul risque}

\begin{problem}[{Problème du week-end --- deux portefeuilles, un seul risque}]\label{pb:hf:extensions-of-the-inventory-framework:1}
Un desk fait du market making sur deux instruments qui évoluent ensemble et dispose d'un signal pour l'un d'entre eux.

\medskip\noindent\textbf{Partie I --- Bornes et horizon long.}
\begin{enumerate}
\item Définir une \omterm{def:hf:extensions-of-the-inventory-framework:bound}{borne d'inventaire} et dire pourquoi elle ramène les cotations à long horizon à un
problème de vecteur propre.
\item Démontrer la \cref{prop:hf:extensions-of-the-inventory-framework:stationary}.
\item Établir la forme fermée de la \cref{prop:hf:extensions-of-the-inventory-framework:closed}.
\item Préciser quand la forme fermée est une bonne approximation, et donner les pires écarts pour $\phi=0.5$, 5 et 50.
\end{enumerate}

\medskip\noindent\textbf{Partie II --- Deux actifs.}
\begin{enumerate}[resume]
\item Définir le \omterm{def:hf:extensions-of-the-inventory-framework:multi}{market making multi-actifs} et écrire la charge de risque.
\item Expliquer pourquoi les instruments doivent partager $k$ pour la solution exacte.
\item Donner le skew croisé~: la profondeur acheteuse de l'actif 1 en $q_1=0$ lorsqu'on est courte de quatre, à plat ou longue de quatre sur l'actif 2.
\item Décrire la simulation et sa mesure du risque d'inventaire.
\item Énoncer le \emph{résultat nommé}~: la réduction du risque d'inventaire obtenue en cotant les deux actifs comme un seul portefeuille, à gain moyen égal, pour $\phi=0.1$,
et son étendue selon les pénalités.
\item Expliquer pourquoi la réduction de l'écart-type du P\&L est beaucoup plus faible que celle du risque d'inventaire.
\end{enumerate}

\medskip\noindent\textbf{Partie III --- Le signal.}
\begin{enumerate}[resume]
\item Définir la \omterm{def:hf:extensions-of-the-inventory-framework:drift}{cotation ajustée de la dérive} (drift-adjusted quoting).
\item Démontrer la \cref{prop:hf:extensions-of-the-inventory-framework:target}.
\item Donner le gain simulé et la position moyenne pour $\mu=2$, $\phi=0.5$.
\item Évaluer ce que vaut la dérive si la \omterm{def:hf:extensions-of-the-inventory-framework:bound}{borne d'inventaire} du teneur de marché est inférieure à $q_0$.
\end{enumerate}

\medskip\noindent\textbf{Partie IV --- Sélection adverse.}
\begin{enumerate}[resume]
\item Démontrer les deux parties de la \cref{prop:hf:extensions-of-the-inventory-framework:adverse}.
\item Donner le profit simulé et l'objectif pénalisé sans élargissement, pour $\xi/2$ et pour $\xi$.
\item Expliquer pourquoi le pas de temps grossier orientait vers l'absence d'élargissement.
\item Indiquer laquelle des quatre extensions du chapitre il faudrait implémenter en premier dans un quoteur de production, et pourquoi.
\item Expliquer comment estimer $\xi$ à partir de ses propres exécutions.
\item Énoncer en une phrase ce que les bornes, la corrélation, les signaux et la sélection adverse ont en commun dans ce cadre.
\end{enumerate}
\end{problem}

\section{Questions d'entretien}

\begin{interviewq}[$\star$ \iqroles{trader}]\label{iq:hf:extensions-of-the-inventory-framework:1}
Vous faites du market making sur un contrat à terme sur indice et sur son fonds indiciel coté (ETF). Vous êtes long du fonds. Comment cela doit-il modifier vos cotations sur le contrat à terme~?
\end{interviewq}

\begin{interviewq}[$\star\star$ \iqroles{chercheur}]\label{iq:hf:extensions-of-the-inventory-framework:2}
Pourquoi les cotations à horizon fini cessent-elles de dépendre du temps loin de l'horizon~?
\end{interviewq}

\begin{interviewq}[$\star\star$ \iqroles{chercheur}]\label{iq:hf:extensions-of-the-inventory-framework:3}
Vos exécutions sont suivies en moyenne d'un mouvement adverse de 0.3 tick. De combien élargissez-vous, et que devez-vous savoir d'abord~?
\end{interviewq}

\begin{interviewq}[$\star\star$ \iqroles{développeur}]\label{iq:hf:extensions-of-the-inventory-framework:4}
Un quoteur doit évaluer ses profondeurs en moins d'une microseconde. Comment livrez-vous un modèle dont la solution exacte est un problème de vecteur propre~?
\end{interviewq}

\begin{interviewq}[$\star\star$ \iqroles{risque}]\label{iq:hf:extensions-of-the-inventory-framework:5}
Un desk mesure le risque d'inventaire instrument par instrument. Qu'est-ce qui lui échappe, et comment le mesureriez-vous à la place~?
\end{interviewq}

\begin{interviewq}[$\star\star\star$ \iqroles{chercheur}]\label{iq:hf:extensions-of-the-inventory-framework:6}
Dans un modèle de market making à pénalité continue, montrer qu'une dérive constante $\mu$ conduit le teneur de marché à viser la position $\mu/(2\phi)$.
\end{interviewq}
